% CNETA logo -- TikZ source for every exported variant.
%
% The mark is a rooted phylogeny whose four tips run into copy-number
% profiles, drawn as heatmap rows the way copy-number trees are usually
% plotted. Clades share events: the gain in segment 2 is inherited by the
% top clade, the loss in segment 1 by its inner cherry.
%
% Plain pdflatex on TeX Live "basic" (TikZ + Latin Modern, nothing else).
% Do not run this file directly; ./build.sh renders all variants. To render
% one by hand:
%
%   pdflatex -jobname=cneta-logo-dark \
%     "\def\LogoLayout{full}\def\LogoScheme{dark}% CNETA logo -- TikZ source for every exported variant.
%
% The mark is a rooted phylogeny whose four tips run into copy-number
% profiles, drawn as heatmap rows the way copy-number trees are usually
% plotted. Clades share events: the gain in segment 2 is inherited by the
% top clade, the loss in segment 1 by its inner cherry.
%
% Plain pdflatex on TeX Live "basic" (TikZ + Latin Modern, nothing else).
% Do not run this file directly; ./build.sh renders all variants. To render
% one by hand:
%
%   pdflatex -jobname=cneta-logo-dark \
%     "\def\LogoLayout{full}\def\LogoScheme{dark}% CNETA logo -- TikZ source for every exported variant.
%
% The mark is a rooted phylogeny whose four tips run into copy-number
% profiles, drawn as heatmap rows the way copy-number trees are usually
% plotted. Clades share events: the gain in segment 2 is inherited by the
% top clade, the loss in segment 1 by its inner cherry.
%
% Plain pdflatex on TeX Live "basic" (TikZ + Latin Modern, nothing else).
% Do not run this file directly; ./build.sh renders all variants. To render
% one by hand:
%
%   pdflatex -jobname=cneta-logo-dark \
%     "\def\LogoLayout{full}\def\LogoScheme{dark}% CNETA logo -- TikZ source for every exported variant.
%
% The mark is a rooted phylogeny whose four tips run into copy-number
% profiles, drawn as heatmap rows the way copy-number trees are usually
% plotted. Clades share events: the gain in segment 2 is inherited by the
% top clade, the loss in segment 1 by its inner cherry.
%
% Plain pdflatex on TeX Live "basic" (TikZ + Latin Modern, nothing else).
% Do not run this file directly; ./build.sh renders all variants. To render
% one by hand:
%
%   pdflatex -jobname=cneta-logo-dark \
%     "\def\LogoLayout{full}\def\LogoScheme{dark}\input{cneta-logo}"
%
%   \LogoLayout  full | tagline | icon | favicon   (default full)
%   \LogoScheme  light | dark                      (default light) ink colours
%   \LogoBg      none | light | dark               (default none = transparent)
%
% `favicon` is a simplified mark (3 tips, chunkier strokes) on a dark tile,
% so that it still reads at 16 px; it ignores \LogoScheme and \LogoBg.

\providecommand\LogoLayout{full}
\providecommand\LogoScheme{light}
\providecommand\LogoBg{none}

\documentclass{article}
\usepackage[T1]{fontenc}
\usepackage{lmodern}
\usepackage{tikz}
\usetikzlibrary{calc}
\pagestyle{empty}

% ---------------------------------------------------------------- palette
% Copy-number states: L = loss, N = neutral, G = gain, A = amplification.
\definecolor{cnInkLight}{HTML}{1D2B36}   % tree + wordmark on light ground
\definecolor{cnInkDark}{HTML}{FFFFFF}    % ... on dark ground
\definecolor{cnLoss}{HTML}{3B82C4}
\definecolor{cnGain}{HTML}{F08A3C}
\definecolor{cnAmp}{HTML}{D1382E}
\definecolor{cnNeutralLight}{HTML}{CBD5DE}
\definecolor{cnNeutralDark}{HTML}{5B6E7E}
\definecolor{cnBgLight}{HTML}{FFFFFF}
\definecolor{cnBgDark}{HTML}{1D2B36}

\def\tmp{dark}
\ifx\LogoScheme\tmp
  \colorlet{cnInk}{cnInkDark}\colorlet{cnNeutral}{cnNeutralDark}
\else
  \colorlet{cnInk}{cnInkLight}\colorlet{cnNeutral}{cnNeutralLight}
\fi
\colorlet{cnL}{cnLoss}\colorlet{cnN}{cnNeutral}
\colorlet{cnG}{cnGain}\colorlet{cnA}{cnAmp}

% ------------------------------------------------------------------ mark
% Drawn in a 10 x 10 box; 1 unit = 1 mm.
\tikzset{cntree/.style={draw=cnInk, line width=0.78mm,
                        line cap=round, line join=round}}

\def\tipA{8.3}\def\tipB{6.1}\def\tipC{3.9}\def\tipD{1.7}
\def\rowH{1.62}                % height of a copy-number row
\def\rowX{6.25}                % left edge of the rows
\def\segW{0.94}                % width of one genome segment (4 per row)

% \cnrow{y}{s1,s2,s3,s4}: one copy-number profile, states L/N/G/A
\newcommand\cnrow[2]{%
  \begin{scope}
    \clip[rounded corners=0.32mm] (\rowX,#1-\rowH/2)
      rectangle ++(4*\segW,\rowH);
    \foreach \s [count=\i from 0] in {#2}
      \fill[cn\s] (\rowX+\i*\segW,#1-\rowH/2) rectangle ++(\segW+0.01,\rowH);
  \end{scope}}

\newcommand\cnicon{%
  % tree: ((A,(B,C)),D), rectangular, root on the left
  \pgfmathsetmacro\yBC{(\tipB+\tipC)/2}
  \pgfmathsetmacro\yX{(\tipA+\yBC)/2}
  \pgfmathsetmacro\yR{(\yX+\tipD)/2}
  \begin{scope}[cntree]
    \draw (0.4,\yR) -- (1.5,\yR);                       % root stem
    \draw (1.5,\tipD) |- (2.9,\yX) (1.5,\tipD) -- (5.65,\tipD);
    \draw (2.9,\tipA) -- (2.9,\yBC) (2.9,\tipA) -- (5.65,\tipA);
    \draw (2.9,\yBC) -- (4.3,\yBC);
    \draw (4.3,\tipB) -- (4.3,\tipC)
          (4.3,\tipB) -- (5.65,\tipB) (4.3,\tipC) -- (5.65,\tipC);
  \end{scope}
  % copy-number profiles at the tips
  \cnrow{\tipA}{N,G,N,L}
  \cnrow{\tipB}{L,G,A,N}
  \cnrow{\tipC}{L,G,N,G}
  \cnrow{\tipD}{G,N,L,N}}

% Favicon mark: ((A,B),C) with two segments per row, always light-on-dark.
% Strokes are ~1.7 px with ~1.8 px gaps and rows ~3.5 px tall at 16 px, about
% the least that still reads as "a tree ending in coloured blocks".
\newcommand\cnfavicon{%
  \colorlet{cnInk}{cnInkDark}\colorlet{cnN}{cnNeutralLight}
  \begin{scope}[cntree, line width=1.0mm]
    \draw (1.2,4.35) -- (2.2,4.35);                     % root stem
    \draw (2.2,2.3) |- (4.3,6.4);                       % C, then the cherry
    \draw (2.2,2.3) -- (6.0,2.3);
    \draw (4.3,5.1) -- (4.3,7.7) (4.3,5.1) -- (6.0,5.1) (4.3,7.7) -- (6.0,7.7);
  \end{scope}
  \foreach \y/\a/\b in {7.7/G/N, 5.1/G/A, 2.3/N/L} {
    \begin{scope}
      \clip[rounded corners=0.3mm] (6.0,\y-0.95) rectangle ++(3.2,1.9);
      \fill[cn\a] (6.0,\y-0.95) rectangle ++(1.6,1.9);
      \fill[cn\b] (7.6,\y-0.95) rectangle ++(1.6,1.9);
    \end{scope}}}

% ---------------------------------------------------------------- output
% Ship the picture as a page cropped exactly to its bounding box.
\newcommand\shippicture[1]{%
  \setbox0\hbox{#1}%
  \pdfpagewidth=\wd0 \pdfpageheight=\dimexpr\ht0+\dp0\relax
  \hoffset=-1in \voffset=-1in
  \shipout\vbox{\box0}}

\pgfdeclarelayer{background}\pgfsetlayers{background,main}
\def\tmpnone{none}\def\tmpdark{dark}
% Padding around the drawing; the optional background tile fills it.
\def\pad{0.9}

\begin{document}
\shippicture{%
\begin{tikzpicture}[x=1mm, y=1mm]
  \def\tmp{favicon}
  \ifx\LogoLayout\tmp
    \def\pad{0}\def\LogoBg{dark}
    \cnfavicon
  \else
    \cnicon
  \fi
  \def\tmp{full}
  \ifx\LogoLayout\tmp
    \node[anchor=base west, inner sep=0, text=cnInk,
          font=\fontsize{23}{23}\sffamily\bfseries]
      at (12.4,2.2) {CNETA};
  \fi
  \def\tmp{tagline}
  \ifx\LogoLayout\tmp
    \node[anchor=base west, inner sep=0, text=cnInk,
          font=\fontsize{21}{21}\sffamily\bfseries]
      at (12.4,4.1) {CNETA};
    \node[anchor=base west, inner sep=0, text=cnInk,
          font=\fontsize{4.1}{5}\sffamily]
      at (12.55,0.75) {copy number evolutionary tree analysis};
  \fi
  % Fixed extent: the icon's 10 mm box plus padding, so layouts align.
  \path (0,0) (10,10);
  \path let \p1=(current bounding box.north east) in
    coordinate (bbNE) at (\x1+\pad mm, 10mm+\pad mm)
    coordinate (bbSW) at (-\pad mm, -\pad mm);
  \useasboundingbox (bbSW) rectangle (bbNE);
  \ifx\LogoBg\tmpnone\else
    \ifx\LogoBg\tmpdark\colorlet{cnBg}{cnBgDark}\else\colorlet{cnBg}{cnBgLight}\fi
    \begin{pgfonlayer}{background}
      \fill[cnBg, rounded corners=1.8mm] (bbSW) rectangle (bbNE);
    \end{pgfonlayer}
  \fi
\end{tikzpicture}}
\end{document}
"
%
%   \LogoLayout  full | tagline | icon | favicon   (default full)
%   \LogoScheme  light | dark                      (default light) ink colours
%   \LogoBg      none | light | dark               (default none = transparent)
%
% `favicon` is a simplified mark (3 tips, chunkier strokes) on a dark tile,
% so that it still reads at 16 px; it ignores \LogoScheme and \LogoBg.

\providecommand\LogoLayout{full}
\providecommand\LogoScheme{light}
\providecommand\LogoBg{none}

\documentclass{article}
\usepackage[T1]{fontenc}
\usepackage{lmodern}
\usepackage{tikz}
\usetikzlibrary{calc}
\pagestyle{empty}

% ---------------------------------------------------------------- palette
% Copy-number states: L = loss, N = neutral, G = gain, A = amplification.
\definecolor{cnInkLight}{HTML}{1D2B36}   % tree + wordmark on light ground
\definecolor{cnInkDark}{HTML}{FFFFFF}    % ... on dark ground
\definecolor{cnLoss}{HTML}{3B82C4}
\definecolor{cnGain}{HTML}{F08A3C}
\definecolor{cnAmp}{HTML}{D1382E}
\definecolor{cnNeutralLight}{HTML}{CBD5DE}
\definecolor{cnNeutralDark}{HTML}{5B6E7E}
\definecolor{cnBgLight}{HTML}{FFFFFF}
\definecolor{cnBgDark}{HTML}{1D2B36}

\def\tmp{dark}
\ifx\LogoScheme\tmp
  \colorlet{cnInk}{cnInkDark}\colorlet{cnNeutral}{cnNeutralDark}
\else
  \colorlet{cnInk}{cnInkLight}\colorlet{cnNeutral}{cnNeutralLight}
\fi
\colorlet{cnL}{cnLoss}\colorlet{cnN}{cnNeutral}
\colorlet{cnG}{cnGain}\colorlet{cnA}{cnAmp}

% ------------------------------------------------------------------ mark
% Drawn in a 10 x 10 box; 1 unit = 1 mm.
\tikzset{cntree/.style={draw=cnInk, line width=0.78mm,
                        line cap=round, line join=round}}

\def\tipA{8.3}\def\tipB{6.1}\def\tipC{3.9}\def\tipD{1.7}
\def\rowH{1.62}                % height of a copy-number row
\def\rowX{6.25}                % left edge of the rows
\def\segW{0.94}                % width of one genome segment (4 per row)

% \cnrow{y}{s1,s2,s3,s4}: one copy-number profile, states L/N/G/A
\newcommand\cnrow[2]{%
  \begin{scope}
    \clip[rounded corners=0.32mm] (\rowX,#1-\rowH/2)
      rectangle ++(4*\segW,\rowH);
    \foreach \s [count=\i from 0] in {#2}
      \fill[cn\s] (\rowX+\i*\segW,#1-\rowH/2) rectangle ++(\segW+0.01,\rowH);
  \end{scope}}

\newcommand\cnicon{%
  % tree: ((A,(B,C)),D), rectangular, root on the left
  \pgfmathsetmacro\yBC{(\tipB+\tipC)/2}
  \pgfmathsetmacro\yX{(\tipA+\yBC)/2}
  \pgfmathsetmacro\yR{(\yX+\tipD)/2}
  \begin{scope}[cntree]
    \draw (0.4,\yR) -- (1.5,\yR);                       % root stem
    \draw (1.5,\tipD) |- (2.9,\yX) (1.5,\tipD) -- (5.65,\tipD);
    \draw (2.9,\tipA) -- (2.9,\yBC) (2.9,\tipA) -- (5.65,\tipA);
    \draw (2.9,\yBC) -- (4.3,\yBC);
    \draw (4.3,\tipB) -- (4.3,\tipC)
          (4.3,\tipB) -- (5.65,\tipB) (4.3,\tipC) -- (5.65,\tipC);
  \end{scope}
  % copy-number profiles at the tips
  \cnrow{\tipA}{N,G,N,L}
  \cnrow{\tipB}{L,G,A,N}
  \cnrow{\tipC}{L,G,N,G}
  \cnrow{\tipD}{G,N,L,N}}

% Favicon mark: ((A,B),C) with two segments per row, always light-on-dark.
% Strokes are ~1.7 px with ~1.8 px gaps and rows ~3.5 px tall at 16 px, about
% the least that still reads as "a tree ending in coloured blocks".
\newcommand\cnfavicon{%
  \colorlet{cnInk}{cnInkDark}\colorlet{cnN}{cnNeutralLight}
  \begin{scope}[cntree, line width=1.0mm]
    \draw (1.2,4.35) -- (2.2,4.35);                     % root stem
    \draw (2.2,2.3) |- (4.3,6.4);                       % C, then the cherry
    \draw (2.2,2.3) -- (6.0,2.3);
    \draw (4.3,5.1) -- (4.3,7.7) (4.3,5.1) -- (6.0,5.1) (4.3,7.7) -- (6.0,7.7);
  \end{scope}
  \foreach \y/\a/\b in {7.7/G/N, 5.1/G/A, 2.3/N/L} {
    \begin{scope}
      \clip[rounded corners=0.3mm] (6.0,\y-0.95) rectangle ++(3.2,1.9);
      \fill[cn\a] (6.0,\y-0.95) rectangle ++(1.6,1.9);
      \fill[cn\b] (7.6,\y-0.95) rectangle ++(1.6,1.9);
    \end{scope}}}

% ---------------------------------------------------------------- output
% Ship the picture as a page cropped exactly to its bounding box.
\newcommand\shippicture[1]{%
  \setbox0\hbox{#1}%
  \pdfpagewidth=\wd0 \pdfpageheight=\dimexpr\ht0+\dp0\relax
  \hoffset=-1in \voffset=-1in
  \shipout\vbox{\box0}}

\pgfdeclarelayer{background}\pgfsetlayers{background,main}
\def\tmpnone{none}\def\tmpdark{dark}
% Padding around the drawing; the optional background tile fills it.
\def\pad{0.9}

\begin{document}
\shippicture{%
\begin{tikzpicture}[x=1mm, y=1mm]
  \def\tmp{favicon}
  \ifx\LogoLayout\tmp
    \def\pad{0}\def\LogoBg{dark}
    \cnfavicon
  \else
    \cnicon
  \fi
  \def\tmp{full}
  \ifx\LogoLayout\tmp
    \node[anchor=base west, inner sep=0, text=cnInk,
          font=\fontsize{23}{23}\sffamily\bfseries]
      at (12.4,2.2) {CNETA};
  \fi
  \def\tmp{tagline}
  \ifx\LogoLayout\tmp
    \node[anchor=base west, inner sep=0, text=cnInk,
          font=\fontsize{21}{21}\sffamily\bfseries]
      at (12.4,4.1) {CNETA};
    \node[anchor=base west, inner sep=0, text=cnInk,
          font=\fontsize{4.1}{5}\sffamily]
      at (12.55,0.75) {copy number evolutionary tree analysis};
  \fi
  % Fixed extent: the icon's 10 mm box plus padding, so layouts align.
  \path (0,0) (10,10);
  \path let \p1=(current bounding box.north east) in
    coordinate (bbNE) at (\x1+\pad mm, 10mm+\pad mm)
    coordinate (bbSW) at (-\pad mm, -\pad mm);
  \useasboundingbox (bbSW) rectangle (bbNE);
  \ifx\LogoBg\tmpnone\else
    \ifx\LogoBg\tmpdark\colorlet{cnBg}{cnBgDark}\else\colorlet{cnBg}{cnBgLight}\fi
    \begin{pgfonlayer}{background}
      \fill[cnBg, rounded corners=1.8mm] (bbSW) rectangle (bbNE);
    \end{pgfonlayer}
  \fi
\end{tikzpicture}}
\end{document}
"
%
%   \LogoLayout  full | tagline | icon | favicon   (default full)
%   \LogoScheme  light | dark                      (default light) ink colours
%   \LogoBg      none | light | dark               (default none = transparent)
%
% `favicon` is a simplified mark (3 tips, chunkier strokes) on a dark tile,
% so that it still reads at 16 px; it ignores \LogoScheme and \LogoBg.

\providecommand\LogoLayout{full}
\providecommand\LogoScheme{light}
\providecommand\LogoBg{none}

\documentclass{article}
\usepackage[T1]{fontenc}
\usepackage{lmodern}
\usepackage{tikz}
\usetikzlibrary{calc}
\pagestyle{empty}

% ---------------------------------------------------------------- palette
% Copy-number states: L = loss, N = neutral, G = gain, A = amplification.
\definecolor{cnInkLight}{HTML}{1D2B36}   % tree + wordmark on light ground
\definecolor{cnInkDark}{HTML}{FFFFFF}    % ... on dark ground
\definecolor{cnLoss}{HTML}{3B82C4}
\definecolor{cnGain}{HTML}{F08A3C}
\definecolor{cnAmp}{HTML}{D1382E}
\definecolor{cnNeutralLight}{HTML}{CBD5DE}
\definecolor{cnNeutralDark}{HTML}{5B6E7E}
\definecolor{cnBgLight}{HTML}{FFFFFF}
\definecolor{cnBgDark}{HTML}{1D2B36}

\def\tmp{dark}
\ifx\LogoScheme\tmp
  \colorlet{cnInk}{cnInkDark}\colorlet{cnNeutral}{cnNeutralDark}
\else
  \colorlet{cnInk}{cnInkLight}\colorlet{cnNeutral}{cnNeutralLight}
\fi
\colorlet{cnL}{cnLoss}\colorlet{cnN}{cnNeutral}
\colorlet{cnG}{cnGain}\colorlet{cnA}{cnAmp}

% ------------------------------------------------------------------ mark
% Drawn in a 10 x 10 box; 1 unit = 1 mm.
\tikzset{cntree/.style={draw=cnInk, line width=0.78mm,
                        line cap=round, line join=round}}

\def\tipA{8.3}\def\tipB{6.1}\def\tipC{3.9}\def\tipD{1.7}
\def\rowH{1.62}                % height of a copy-number row
\def\rowX{6.25}                % left edge of the rows
\def\segW{0.94}                % width of one genome segment (4 per row)

% \cnrow{y}{s1,s2,s3,s4}: one copy-number profile, states L/N/G/A
\newcommand\cnrow[2]{%
  \begin{scope}
    \clip[rounded corners=0.32mm] (\rowX,#1-\rowH/2)
      rectangle ++(4*\segW,\rowH);
    \foreach \s [count=\i from 0] in {#2}
      \fill[cn\s] (\rowX+\i*\segW,#1-\rowH/2) rectangle ++(\segW+0.01,\rowH);
  \end{scope}}

\newcommand\cnicon{%
  % tree: ((A,(B,C)),D), rectangular, root on the left
  \pgfmathsetmacro\yBC{(\tipB+\tipC)/2}
  \pgfmathsetmacro\yX{(\tipA+\yBC)/2}
  \pgfmathsetmacro\yR{(\yX+\tipD)/2}
  \begin{scope}[cntree]
    \draw (0.4,\yR) -- (1.5,\yR);                       % root stem
    \draw (1.5,\tipD) |- (2.9,\yX) (1.5,\tipD) -- (5.65,\tipD);
    \draw (2.9,\tipA) -- (2.9,\yBC) (2.9,\tipA) -- (5.65,\tipA);
    \draw (2.9,\yBC) -- (4.3,\yBC);
    \draw (4.3,\tipB) -- (4.3,\tipC)
          (4.3,\tipB) -- (5.65,\tipB) (4.3,\tipC) -- (5.65,\tipC);
  \end{scope}
  % copy-number profiles at the tips
  \cnrow{\tipA}{N,G,N,L}
  \cnrow{\tipB}{L,G,A,N}
  \cnrow{\tipC}{L,G,N,G}
  \cnrow{\tipD}{G,N,L,N}}

% Favicon mark: ((A,B),C) with two segments per row, always light-on-dark.
% Strokes are ~1.7 px with ~1.8 px gaps and rows ~3.5 px tall at 16 px, about
% the least that still reads as "a tree ending in coloured blocks".
\newcommand\cnfavicon{%
  \colorlet{cnInk}{cnInkDark}\colorlet{cnN}{cnNeutralLight}
  \begin{scope}[cntree, line width=1.0mm]
    \draw (1.2,4.35) -- (2.2,4.35);                     % root stem
    \draw (2.2,2.3) |- (4.3,6.4);                       % C, then the cherry
    \draw (2.2,2.3) -- (6.0,2.3);
    \draw (4.3,5.1) -- (4.3,7.7) (4.3,5.1) -- (6.0,5.1) (4.3,7.7) -- (6.0,7.7);
  \end{scope}
  \foreach \y/\a/\b in {7.7/G/N, 5.1/G/A, 2.3/N/L} {
    \begin{scope}
      \clip[rounded corners=0.3mm] (6.0,\y-0.95) rectangle ++(3.2,1.9);
      \fill[cn\a] (6.0,\y-0.95) rectangle ++(1.6,1.9);
      \fill[cn\b] (7.6,\y-0.95) rectangle ++(1.6,1.9);
    \end{scope}}}

% ---------------------------------------------------------------- output
% Ship the picture as a page cropped exactly to its bounding box.
\newcommand\shippicture[1]{%
  \setbox0\hbox{#1}%
  \pdfpagewidth=\wd0 \pdfpageheight=\dimexpr\ht0+\dp0\relax
  \hoffset=-1in \voffset=-1in
  \shipout\vbox{\box0}}

\pgfdeclarelayer{background}\pgfsetlayers{background,main}
\def\tmpnone{none}\def\tmpdark{dark}
% Padding around the drawing; the optional background tile fills it.
\def\pad{0.9}

\begin{document}
\shippicture{%
\begin{tikzpicture}[x=1mm, y=1mm]
  \def\tmp{favicon}
  \ifx\LogoLayout\tmp
    \def\pad{0}\def\LogoBg{dark}
    \cnfavicon
  \else
    \cnicon
  \fi
  \def\tmp{full}
  \ifx\LogoLayout\tmp
    \node[anchor=base west, inner sep=0, text=cnInk,
          font=\fontsize{23}{23}\sffamily\bfseries]
      at (12.4,2.2) {CNETA};
  \fi
  \def\tmp{tagline}
  \ifx\LogoLayout\tmp
    \node[anchor=base west, inner sep=0, text=cnInk,
          font=\fontsize{21}{21}\sffamily\bfseries]
      at (12.4,4.1) {CNETA};
    \node[anchor=base west, inner sep=0, text=cnInk,
          font=\fontsize{4.1}{5}\sffamily]
      at (12.55,0.75) {copy number evolutionary tree analysis};
  \fi
  % Fixed extent: the icon's 10 mm box plus padding, so layouts align.
  \path (0,0) (10,10);
  \path let \p1=(current bounding box.north east) in
    coordinate (bbNE) at (\x1+\pad mm, 10mm+\pad mm)
    coordinate (bbSW) at (-\pad mm, -\pad mm);
  \useasboundingbox (bbSW) rectangle (bbNE);
  \ifx\LogoBg\tmpnone\else
    \ifx\LogoBg\tmpdark\colorlet{cnBg}{cnBgDark}\else\colorlet{cnBg}{cnBgLight}\fi
    \begin{pgfonlayer}{background}
      \fill[cnBg, rounded corners=1.8mm] (bbSW) rectangle (bbNE);
    \end{pgfonlayer}
  \fi
\end{tikzpicture}}
\end{document}
"
%
%   \LogoLayout  full | tagline | icon | favicon   (default full)
%   \LogoScheme  light | dark                      (default light) ink colours
%   \LogoBg      none | light | dark               (default none = transparent)
%
% `favicon` is a simplified mark (3 tips, chunkier strokes) on a dark tile,
% so that it still reads at 16 px; it ignores \LogoScheme and \LogoBg.

\providecommand\LogoLayout{full}
\providecommand\LogoScheme{light}
\providecommand\LogoBg{none}

\documentclass{article}
\usepackage[T1]{fontenc}
\usepackage{lmodern}
\usepackage{tikz}
\usetikzlibrary{calc}
\pagestyle{empty}

% ---------------------------------------------------------------- palette
% Copy-number states: L = loss, N = neutral, G = gain, A = amplification.
\definecolor{cnInkLight}{HTML}{1D2B36}   % tree + wordmark on light ground
\definecolor{cnInkDark}{HTML}{FFFFFF}    % ... on dark ground
\definecolor{cnLoss}{HTML}{3B82C4}
\definecolor{cnGain}{HTML}{F08A3C}
\definecolor{cnAmp}{HTML}{D1382E}
\definecolor{cnNeutralLight}{HTML}{CBD5DE}
\definecolor{cnNeutralDark}{HTML}{5B6E7E}
\definecolor{cnBgLight}{HTML}{FFFFFF}
\definecolor{cnBgDark}{HTML}{1D2B36}

\def\tmp{dark}
\ifx\LogoScheme\tmp
  \colorlet{cnInk}{cnInkDark}\colorlet{cnNeutral}{cnNeutralDark}
\else
  \colorlet{cnInk}{cnInkLight}\colorlet{cnNeutral}{cnNeutralLight}
\fi
\colorlet{cnL}{cnLoss}\colorlet{cnN}{cnNeutral}
\colorlet{cnG}{cnGain}\colorlet{cnA}{cnAmp}

% ------------------------------------------------------------------ mark
% Drawn in a 10 x 10 box; 1 unit = 1 mm.
\tikzset{cntree/.style={draw=cnInk, line width=0.78mm,
                        line cap=round, line join=round}}

\def\tipA{8.3}\def\tipB{6.1}\def\tipC{3.9}\def\tipD{1.7}
\def\rowH{1.62}                % height of a copy-number row
\def\rowX{6.25}                % left edge of the rows
\def\segW{0.94}                % width of one genome segment (4 per row)

% \cnrow{y}{s1,s2,s3,s4}: one copy-number profile, states L/N/G/A
\newcommand\cnrow[2]{%
  \begin{scope}
    \clip[rounded corners=0.32mm] (\rowX,#1-\rowH/2)
      rectangle ++(4*\segW,\rowH);
    \foreach \s [count=\i from 0] in {#2}
      \fill[cn\s] (\rowX+\i*\segW,#1-\rowH/2) rectangle ++(\segW+0.01,\rowH);
  \end{scope}}

\newcommand\cnicon{%
  % tree: ((A,(B,C)),D), rectangular, root on the left
  \pgfmathsetmacro\yBC{(\tipB+\tipC)/2}
  \pgfmathsetmacro\yX{(\tipA+\yBC)/2}
  \pgfmathsetmacro\yR{(\yX+\tipD)/2}
  \begin{scope}[cntree]
    \draw (0.4,\yR) -- (1.5,\yR);                       % root stem
    \draw (1.5,\tipD) |- (2.9,\yX) (1.5,\tipD) -- (5.65,\tipD);
    \draw (2.9,\tipA) -- (2.9,\yBC) (2.9,\tipA) -- (5.65,\tipA);
    \draw (2.9,\yBC) -- (4.3,\yBC);
    \draw (4.3,\tipB) -- (4.3,\tipC)
          (4.3,\tipB) -- (5.65,\tipB) (4.3,\tipC) -- (5.65,\tipC);
  \end{scope}
  % copy-number profiles at the tips
  \cnrow{\tipA}{N,G,N,L}
  \cnrow{\tipB}{L,G,A,N}
  \cnrow{\tipC}{L,G,N,G}
  \cnrow{\tipD}{G,N,L,N}}

% Favicon mark: ((A,B),C) with two segments per row, always light-on-dark.
% Strokes are ~1.7 px with ~1.8 px gaps and rows ~3.5 px tall at 16 px, about
% the least that still reads as "a tree ending in coloured blocks".
\newcommand\cnfavicon{%
  \colorlet{cnInk}{cnInkDark}\colorlet{cnN}{cnNeutralLight}
  \begin{scope}[cntree, line width=1.0mm]
    \draw (1.2,4.35) -- (2.2,4.35);                     % root stem
    \draw (2.2,2.3) |- (4.3,6.4);                       % C, then the cherry
    \draw (2.2,2.3) -- (6.0,2.3);
    \draw (4.3,5.1) -- (4.3,7.7) (4.3,5.1) -- (6.0,5.1) (4.3,7.7) -- (6.0,7.7);
  \end{scope}
  \foreach \y/\a/\b in {7.7/G/N, 5.1/G/A, 2.3/N/L} {
    \begin{scope}
      \clip[rounded corners=0.3mm] (6.0,\y-0.95) rectangle ++(3.2,1.9);
      \fill[cn\a] (6.0,\y-0.95) rectangle ++(1.6,1.9);
      \fill[cn\b] (7.6,\y-0.95) rectangle ++(1.6,1.9);
    \end{scope}}}

% ---------------------------------------------------------------- output
% Ship the picture as a page cropped exactly to its bounding box.
\newcommand\shippicture[1]{%
  \setbox0\hbox{#1}%
  \pdfpagewidth=\wd0 \pdfpageheight=\dimexpr\ht0+\dp0\relax
  \hoffset=-1in \voffset=-1in
  \shipout\vbox{\box0}}

\pgfdeclarelayer{background}\pgfsetlayers{background,main}
\def\tmpnone{none}\def\tmpdark{dark}
% Padding around the drawing; the optional background tile fills it.
\def\pad{0.9}

\begin{document}
\shippicture{%
\begin{tikzpicture}[x=1mm, y=1mm]
  \def\tmp{favicon}
  \ifx\LogoLayout\tmp
    \def\pad{0}\def\LogoBg{dark}
    \cnfavicon
  \else
    \cnicon
  \fi
  \def\tmp{full}
  \ifx\LogoLayout\tmp
    \node[anchor=base west, inner sep=0, text=cnInk,
          font=\fontsize{23}{23}\sffamily\bfseries]
      at (12.4,2.2) {CNETA};
  \fi
  \def\tmp{tagline}
  \ifx\LogoLayout\tmp
    \node[anchor=base west, inner sep=0, text=cnInk,
          font=\fontsize{21}{21}\sffamily\bfseries]
      at (12.4,4.1) {CNETA};
    \node[anchor=base west, inner sep=0, text=cnInk,
          font=\fontsize{4.1}{5}\sffamily]
      at (12.55,0.75) {copy number evolutionary tree analysis};
  \fi
  % Fixed extent: the icon's 10 mm box plus padding, so layouts align.
  \path (0,0) (10,10);
  \path let \p1=(current bounding box.north east) in
    coordinate (bbNE) at (\x1+\pad mm, 10mm+\pad mm)
    coordinate (bbSW) at (-\pad mm, -\pad mm);
  \useasboundingbox (bbSW) rectangle (bbNE);
  \ifx\LogoBg\tmpnone\else
    \ifx\LogoBg\tmpdark\colorlet{cnBg}{cnBgDark}\else\colorlet{cnBg}{cnBgLight}\fi
    \begin{pgfonlayer}{background}
      \fill[cnBg, rounded corners=1.8mm] (bbSW) rectangle (bbNE);
    \end{pgfonlayer}
  \fi
\end{tikzpicture}}
\end{document}
